\section{Analytic preliminaries}
\label{sec:prelim}
We record the analytic inputs used in the proof, including the parameter
estimates needed for moving balls. Sobolev spaces are defined in smooth
charts by $\|v\|_{H^s}^2=\int\langle\xi\rangle^{2s}|\widehat v(\xi)|^2\,d\xi$,
using $\widehat v=(2\pi)^{-n/2}\int e^{-ix\cdot\xi}v(x)\,dx$.
A finite partition of unity defines the norms on compact manifolds.
Different smooth charts and positive smooth densities give equivalent
norms; all such comparisons below are on fixed compact sets.
The measure-theoretic inputs are monotone and dominated convergence,
Tonelli and Fubini on sigma-finite spaces, and completeness of $L^2$;
see \citep[Theorems~3.11, 3.31, 5.28, 5.32 and~7.24]{Axler2020}.
All measures used here are positive smooth densities, Euclidean
Lebesgue measure or counting measure, hence sigma-finite. Fubini is
used with an absolute-integrability bound; nonnegative square integrals
use Tonelli. Complex formulas follow by real and imaginary parts.
Sobolev spaces are Hilbert completions, and local PDE regularity fixes
their smooth representatives when point values are used.

\subsection{Dirichlet solutions and uniform estimates}
The Fourier inversion and Plancherel statements used here are
\citep[Propositions~3.1--3.2]{TaylorFourier}; the corresponding Fourier
series identity is given there in Section~1.
For a half-space, the operator
\[
 Ef(x',s)=\mathcal F^{-1}_{\xi\to x'}
           \big(e^{-s\langle\xi\rangle}\widehat f(\xi)\big),\qquad s\ge0,
\]
followed by a fixed cutoff, is a bounded right inverse of the trace
$H^1\to H^{1/2}$. Indeed integration in $s$ gives
$\|Ef\|_{H^1}^2\le C\int\langle\xi\rangle|\widehat f|^2$.
Conversely the one-dimensional fundamental theorem of calculus gives
\[
 \langle\xi\rangle|\widehat u(\xi,0)|^2
 \le\int_0^\infty
       \big(\langle\xi\rangle^2|\widehat u|^2
                        +|\partial_s\widehat u|^2\big)\,ds
\]
for smooth functions of compact support. Density and boundary charts
give the trace and extension assertions on $M$. Convolution after
local extension gives smooth approximation. Zero extension across a
flat boundary with zero trace has no distributional jump; conversely
translations followed by convolution approximate such extensions by
functions supported on the interior side. Applying cutoffs proves
$\ker\operatorname{Tr}=H^1_0(M)$ and the gluing statement for $H^1$
functions with equal traces on a smooth interface.

For completeness the coercivity used here is also present in
\citep[Section~1, (1.16)]{TaylorElliptic}: for a fixed trace $g$ there
are $B>0,C_g$ with
$\frac12\|du\|_2^2\ge B\|u\|_2^2-C_g$ on that trace class.
With $g=0$, replace $u$ by $a u$ and let $a\to\infty$.
Thus $\|u\|_2\le C\|du\|_2$ on $H^1_0(M)$.
The form $a_g(u,v)=\int\langle du,d\bar v\rangle_gdV_g$
therefore defines an equivalent Hilbert norm. The Riesz representation
theorem \citep[Proposition~2.2]{TaylorFunctional} gives a unique zero-trace
solution of $L_gu=F$ for every continuous functional $F$ on $H^1_0$.
Applying it to $-a_g(Ef,\cdot)$ gives the harmonic extension of $f$.
Its energy is at most that of $Ef$, and hence $Q_g$ is bounded on
$H^{1/2}\times H^{1/2}$. Symmetry of $a_g$ gives a self-adjoint
Dirichlet inverse on smooth sources.

Here is the quantitative regularity fact used below. For a smooth
uniformly positive metric on a fixed smooth compact domain and
integer $q\ge0$, its zero-Dirichlet solution obeys
\begin{equation}\label{prelim:dir-estimate}
 \|u\|_{H^{q+2}}\le C_q(\|L_gu\|_{H^q}+\|u\|_{H^1}).
\end{equation}
The constants depend on finitely many coefficient bounds, ellipticity,
and fixed boundary charts. One can derive this estimate directly as
follows. In a flattened boundary chart the weak equation is in
divergence form with symmetric coefficients $a^{ij}$ and $a^{nn}>c>0$.
For a tangential difference quotient $D_{h,k}$, test the equation with
$-D_{-h,k}(\eta^2D_{h,k}u)$. Integration of difference quotients moves
$D_{h,k}$ to $a^{ij}\partial_j u$; its principal term is
\[
 \int\eta^2a^{ij}(x+he_k)
          \partial_jD_{h,k}u\,\partial_iD_{h,k}\bar u.
\]
The coefficient difference and cutoff terms are bounded, by Cauchy--Schwarz
and Young's inequality, by half that positive term plus
$C(\|u\|_{H^1}^2+\|L_gu\|_2^2)$. The source term uses
$\|D_{-h,k}(\eta^2D_{h,k}u)\|_2
 \le C\|\nabla(\eta^2D_{h,k}u)\|_2$.
Passing to $h\to0$ controls every second derivative with at least one
tangential index. The equation solves for $\partial_n^2u$, so $u\in H^2$.
For higher orders apply the same argument to each tangential derivative;
its boundary trace is zero. Commutators contain derivatives of coefficients
and lower derivatives of $u$. Recover derivatives with two or more normal
indices successively from
$a^{nn}\partial_n^2u=-L_gu-\sum_{(i,j)\ne(n,n)}a^{ij}\partial_i\partial_j u$
plus the first-order terms. Induction first on total order and then on
normal order proves \eqref{prelim:dir-estimate}. Interior patches use
unrestricted difference quotients. A finite partition and slightly nested
cutoffs give the global estimate. Approximation of the coefficients and
source, or first applying the inequalities to difference quotients of the
weak solution, justifies each step before the asserted derivatives exist.
Nonzero smooth trace is subtracted by a smooth extension. In particular,
smooth data give solutions smooth up to the boundary, in agreement with
\citep[Propositions~1.2--1.3]{TaylorElliptic}.

The same local estimates with cutoffs give
$\|u\|_{H^q(K)}\le C_{q,K}\|u\|_{L^2(D)}$ for homogeneous solutions,
$K\Subset D$, including scalar elliptic equations with bounded smooth
lower-order terms. The initial $H^1$ estimate is obtained by testing with
$\eta^2\bar u$; each further derivative follows from difference quotients.
For distributional solutions local regularity instead follows from the
elliptic parametrix described below.

Pull a smooth family of balls back to the fixed unit ball, including its
limit at radius zero. For each $q$, \eqref{prelim:dir-estimate} and
coercivity give uniform bounded inverses
$L(\theta)^{-1}:H^q\to H^{q+2}\cap H^1_0$ on compact parameter sets.
Multiplying a metric Laplacian by a positive smooth function changes its
source by that function's reciprocal, preserving these assertions.
A perturbation in this operator norm is inverted by the Neumann series.
Differentiation gives
$\partial_\theta R=-R(\partial_\theta L)R$, and its iterates give smooth
parameter dependence in every fixed Sobolev order. Sobolev embedding here
follows directly from Fourier inversion and
$\int\langle\xi\rangle^{-2s}<\infty$ for $s>n/2$.
These estimates control smooth-kernel remainders by taking sufficiently
large, fixed orders; the number of orders may depend on the requested
seminorm, but no shrinking parameter enters that choice.

\subsection{Symbols, Green kernels, and point distributions}
Use classical scalar $\Psi^m=\Psi^m_{1,0}$ operators.
The symbol product, adjoint, $H^s$ boundedness, and elliptic parametrix
are those of \citep[Lecture~I, Sections~2--4 and~8]{TaylorPsido}.
In particular, left symbols satisfy
$b\#c\sim\sum_\alpha(1/i)^{|\alpha|}
 (\partial_\xi^\alpha b)(\partial_x^\alpha c)/\alpha!$;
a scalar commutator lowers the summed order by one, and an operator of
order $m$ maps $H^s$ boundedly to $H^{s-m}$.
An elliptic operator of order two has a two-sided local parametrix
of order $-2$ modulo smoothing. Its kernel is smooth off the diagonal.
For that parametrix, \citep[Lecture~I, (1.23)--(1.24)]{TaylorPsido}
gives $|\partial_x^\alpha\partial_y^\zeta K(x,y)|
 \le C|x-y|^{2-n-|\alpha|-|\zeta|}$ when $n\ge3$.
Subtracting the true Dirichlet inverse leaves a smooth kernel on compact
interior subsets: apply local regularity to the parametrix error, and
to the transposed error in the second variable. Thus these bounds apply
locally to the Dirichlet Green kernel. After dilation, each derivative
has its matching power of the dilation, so compact subsets at fixed
positive separation have uniform bounds. The smooth remainder is harmless
since $\lambda^{n-2}$ and its differentiated factors are bounded for
$n\ge3$. The weak inverse and its transpose supply the Green distribution;
this construction supplies its off-diagonal smooth representative.

The boundary symbol factorization used below is the scalar case of
\citep[Theorem~1.1(i)--(ii) and Proposition~3.1]{JoshiLionheart2005}.
The theorem is stated for an orientable manifold and all form degrees;
we use its local scalar construction in half-balls and normal unit balls,
which are orientable. Its recursion also applies after multiplication
by a positive function: the harmonic solutions are unchanged, and any
chosen transverse derivative trace is a smooth multiple of the unit
normal trace plus a tangential derivative. Its dependence on parameters
is obtained by the symbol recursion and the uniform inverse estimates
above. No global differential-form conclusion is used. The authors'
errata change the sign in a differential-form Green identity and correct
an introductory reference, leaving this scalar construction unchanged.

A point-supported distribution is a finite sum of derivatives of a
point mass \citep[Proposition~4.5]{TaylorFourier}. Locally this also follows
without an additional representation theorem. If it has order $k$ and
all derivatives of a test function through $k$ vanish at the point,
multiplying that function by a cutoff at scale $a$ gives $C^k$ norm
$O(a)$. The distribution's action tends to zero, and equals its action
on the original function by point support. It therefore factors through
the finite $k$-jet space, whose dual is spanned by derivatives of the
point mass. The principal symbol calculation in the next section will
use this fact for the transposed Dirichlet inverse.

\subsection{Continuation and geometric coordinates}
The open-set unique continuation input is
\citep[Theorem~8, Definitions~8.1--8.3]{KochTataru2005} for scalar elliptic
operators, applied in smooth coordinate balls. With positive quadratic
principal symbol, at a nonzero normal covector the simultaneous complex
conditions $p_\theta=\{p_\theta,\theta\}=0$ have no solution:
the imaginary part of the latter bracket is, up to sign,
$2\tau|d\theta|_g^2$. Thus every smooth hypersurface is strongly
pseudoconvex for the surface form of that theorem. One-sided zero sets
propagate across coordinate sphere fronts, and a path with overlapping
balls gives open-set continuation on a connected domain. Distributional
solutions are first made smooth by the elliptic parametrix.
The strict-weight Carleman estimate used for product continuation is
\citep[Definition~8.3 and Theorem~9(a)]{KochTataru2005}; its specialization,
normalization, and uniform margins are also derived in the proof below.

Local existence of geodesics and smooth parameter dependence follow by
writing the geodesic equation as a first-order smooth system, using
\citep[Theorem~2.1, Propositions~6.1--6.2]{TaylorODE}.
The inverse function theorem is \citep[Theorem~3.1]{TaylorODE}.
At zero velocity the exponential map has derivative the identity;
these inputs give normal coordinate balls uniform on compact center sets.
For a compact smooth boundary, the inward normal exponential has
invertible differential at the zero section. Local injectivity and a
finite cover give a uniform collar; otherwise two colliding sequences
would limit to one boundary point and contradict local injectivity there.
Differentiating the geodesic equation gives Jacobi fields and the
$O(r^2)$ center-variation remainder used for the moving traces.
Metric compatibility shows $g(\partial_s,\partial_{z_a})=0$ along a
normal collar: its derivative in $s$ is
$\frac12\partial_{z_a}g(\partial_s,\partial_s)=0$, and it is zero at
$s=0$. Hence $|ds|_g=1$. Any curve reaching the boundary in the collar
has length at least $s$, while the normal segment has length $s$;
curves first leaving a fixed larger collar have a longer fixed lower
bound when $s$ is small. The collar parameter is consequently the
boundary distance. This also proves the boundary-distance smoothness
used in completing the final isometry. Green identities and metric
normal transport use the coordinate divergence formula and integration
by parts against positive densities; orientation is never needed.
